\documentclass[../../../main.tex]{subfiles}
\begin{document}
\subsection{Heisenberg Hamiltonian}
The electron in solid has the Hamiltonian
\begin{equation*}
    H = \sum_{i}\left[\frac{p_i^2}{2m}+V(\mathbf{r}_i)\right]+\frac{1}{2}\sum_{i\neq j}\frac{e^2}{\lvert\mathbf{r}_i-\mathbf{r}_j\rvert}
\end{equation*}
which contains kinetic energy, the periodic crystal potential, and Coulomb repulsion. At this stage there are no explicit spin–spin interactions.
One rewrites this in second quantization and chooses Bloch eigenstates as a basis, which diagonalizes the one-body part. 
Because a full treatment over all bands is impractical, one also projects onto a single partially relevant band near the Fermi level
\begin{equation*}
    H =  \sum_{ij \sigma} \epsilon_{ij} a ^\dagger_{i \sigma} a_{j \sigma}  +
    \frac{1}{2}\sum_{\substack{i j \sigma \\ k l \sigma'}}
    \braket{ i   j  | \frac{1}{r} |   k l   }
    a^\dagger_{ i  \sigma} a^\dagger_{  j  \sigma'} a_{l  \sigma'}a_{  k  \sigma}
\end{equation*}
or further upon decomposition
\begin{multline*}
    H = \sum_{ij\sigma} \epsilon_{ij}\,a_{i\sigma}^{\dagger}a_{j\sigma}
    + \frac{I}{2}\sum_{i\sigma}n_{i\sigma}n_{i-\sigma}
    + \frac{1}{2}\sum_{ij}^{ }{} '  \sum_{\sigma\sigma'}\langle ij|\frac{1}{r}|ij\rangle\,n_{i\sigma}n_{j\sigma'}\\
    -\frac{1}{2}\sum_{ij}{}' \sum_{\sigma\sigma'}J_{ij}a_{i\sigma}^{\dagger}a_{i\sigma'}a_{j\sigma'}^{\dagger}a_{j\sigma}
    + \sum_{ij} {} ' \sum_{\sigma\sigma'}\langle ii|\frac{1}{r}|jj\rangle\,a_{i\sigma}^{\dagger}a_{i\sigma'}^{\dagger}a_{j\sigma'}a_{j\sigma}
\end{multline*}
\begin{enumerate}
    \item \textbf{Hopping term.} It describe electron initially at site $j$ move to  $i$ with amplitude $\epsilon_{ij}$.
    \item \textbf{On-site Coulomb repulsion.} It penalizes double occupation with energy $I$ and leaving single occupied state, where either $n_{i\sigma}$ or $n_{i-\sigma}$ equal to zero.
    \item \textbf{Intersite direct Coulomb repulsion.} ordinary electrostatic repulsion between charge densities on different sites $n_i$ and $n_j$.
    \item \textbf{Exchange interaction.} The origin of Heisenberg Hamiltonian. The spins are exchanged between the two sites: at site $j$ spin $\sigma\to\sigma'$ and at site $i$ $\sigma'\to\sigma$ 
    \item \textbf{Pair hopping.} This moves an electron pair together, two from site $j$ into two $i$.
\end{enumerate}

In the $I\gg \epsilon_{ij}$, double occupancy costs so much energy that low-energy states satisfy $n_i =1$ for all site.
The charge degrees of freedom are then frozen, leaving only local spin variables.
The generate effective Hamiltonian called the Heisenberg Hamiltonian
\begin{equation*}
    H_H = -\sum_{ij } J_{ij} \boldsymbol{\sigma}_i  \cdot \boldsymbol{\sigma}_j
\end{equation*}
or under nearest-neighbor approximation
\begin{equation*}
    H_H = -\sum_{ij } J_{\braket{ij }} \boldsymbol{\sigma}_i  \cdot \boldsymbol{\sigma}_j    
\end{equation*}
Therefore, it is a low-energy effective theory describing interacting localized spins after electronic charge fluctuations have been eliminated.

\subsubsection{Derivation.}
We start from the Hamiltonian
\begin{equation*}
    H = \sum_{i}\left[\frac{p_i^2}{2m}+V(\mathbf{r}_i)\right]+\frac{1}{2}\sum_{i\neq j}\frac{e^2}{\lvert\mathbf{r}_i-\mathbf{r}_j\rvert}
\end{equation*}
The second quantization reads
\begin{equation*}
    H = \sum_{ij}\langle i|h|j\rangle a_{i}^{\dagger}a_{j} +\frac{1}{2}\sum_{ijlm}\langle ij|v|lm\rangle a_{i}^{\dagger}a_{j}^{\dagger}a_{m}a_{l}
\end{equation*}
With the matrix elements
\begin{align*}
    \langle i|h|j\rangle   & = \int d \mathbf{r}\;\phi_{i}^{*}(\mathbf{r}) h(\mathbf{r}) \phi_{j}(\mathbf{r})                                                                              \\
    \langle ij|v|lm\rangle & = \int d \mathbf{r} \;d \mathbf{r}'\; \phi_{i}^{*}(\mathbf{r}) \phi_{j}^{*}(\mathbf{r}') v(\mathbf{r}-\mathbf{r}') \phi_{l}(\mathbf{r}) \phi_{m}(\mathbf{r}')
\end{align*}
Here we use the Bloch state $\ket{n \mathbf{k }\sigma}$.
The translational symmetry of the state implies diagonal one particle operator
\begin{equation*}
    H_1 = \sum_{n\mathbf{k}  \sigma} \braket{n \mathbf{k }\sigma  | h | n \mathbf{k }\sigma} a^\dagger_{n \mathbf{k} \sigma} a_{  n \mathbf{k} \sigma}
    = \sum_{n\mathbf{k}  \sigma} \epsilon_{\mathbf{k }} ^{n} a^\dagger_{n\mathbf{k}  \sigma} a_{n \mathbf{k} \sigma}
\end{equation*}
As for the two particle operator
\begin{align*}
    H_2
     & =   \frac{1}{2}\sum_{\substack{n_1 n_2                                                                                                                             \\ n_1' n_2'}}\sum_{\substack{\mathbf{k}_1 \mathbf{k}_2 \\ \mathbf{k}_1' \mathbf{k}_2'}}\sum_{\substack{\sigma_1 \sigma_2\\\sigma_1' \sigma_2'}}
    \braket{n_1 \mathbf{k}_1 \sigma_1  n_2 \mathbf{k}_2\sigma_2  | \frac{1}{r} |  n_1' \mathbf{k}_1' \sigma_1' n_2'\mathbf{k}_2' \sigma_2'}                               \\
     & \qquad\qquad\times a^\dagger_{n_1\mathbf{k}_1  \sigma_1} a^\dagger_{n_2 \mathbf{k}_2  \sigma_2} a_{ n_2'\mathbf{k}_2'  \sigma_2'}a_{n_1' \mathbf{k}_1'  \sigma_1'} \\
     & =   \frac{1}{2}\sum_{\substack{n_1 n_2                                                                                                                             \\ n_1' n_2'}}\sum_{\substack{\mathbf{k}_1 \mathbf{k}_2 \\ \mathbf{k}_1' \mathbf{k}_2'}}\sum_{\sigma\sigma'}
    \braket{n_1 \mathbf{k}_1  n_2 \mathbf{k}_2  | \frac{1}{r} |  n_1' \mathbf{k}_1' n_2'\mathbf{k}_2'   }
    a^\dagger_{n_1\mathbf{k}_1  \sigma} a^\dagger_{n_2 \mathbf{k}_2  \sigma'} a_{ n_2'\mathbf{k}_2'  \sigma'}a_{n_1' \mathbf{k}_1'  \sigma}
\end{align*}
since the two-body interaction is spin-independent.
The matrix element for the two body operator is
\begin{equation*}
    \langle \mathbf{k}_1 n_1 \mathbf{k}_2 n_2|\frac{1}{r}|\mathbf{k}_1' n_1' \mathbf{k}_2' n_2'\rangle = e^2 \int d\mathbf r\;d\mathbf r'\;
    \frac{\psi_{\mathbf{k}_1}^{n_1}(\mathbf r)\,\psi_{\mathbf{k}_2}^{n_2}(\mathbf r')\,\psi_{\mathbf{k}_1'}^{n_1'}(\mathbf r)\,\psi_{\mathbf{k}_2'}^{n_2'}(\mathbf r')}
    {|\mathbf r-\mathbf r'|}
\end{equation*}
Since it is impractical to handle the whole sum, we perform the single band approximation.
This is only valid if the whole band are energetically well separated, and that the band we retained is the one closest to the Fermi level.
The Hamiltonian now reads
\begin{equation*}
    H =\sum_{\mathbf{k}  \sigma} \epsilon_{\mathbf{k }}  a^\dagger_{\mathbf{k}  \sigma} a_{ \mathbf{k} \sigma}
    +\frac{1}{2}\sum_{\substack{\mathbf{k}_1 \mathbf{k}_2 \sigma \\ \mathbf{k}_1' \mathbf{k}_2' \sigma'}}
    \braket{ \mathbf{k}_1   \mathbf{k}_2  | \frac{1}{r} |   \mathbf{k}_1' \mathbf{k}_2'   }
    a^\dagger_{ \mathbf{k}_1  \sigma} a^\dagger_{  \mathbf{k}_2  \sigma'} a_{\mathbf{k}_2'  \sigma'}a_{  \mathbf{k}_1'  \sigma}
\end{equation*}

We now introduce the Warner function $w(\mathbf{r} - \mathbf{R}_i)$ to transform the Hamiltonian from the momentum space into the real space.
The function is defined as
\begin{equation*}
    \psi_{\mathbf{k}}(\mathbf{r}) = \frac{1}{\sqrt{N}}\sum_{i} w(\mathbf{r}-\mathbf{R}_i)e^{i\mathbf{k}\cdot\mathbf{R}_i}\quad
    w(\mathbf{r}-\mathbf{R}_i) = \frac{1}{\sqrt{N}}\sum_{\mathbf{k}}\psi_{\mathbf{k}}(\mathbf{r})e^{-i\mathbf{k}\cdot\mathbf{R}_i}
\end{equation*}
Note that index $i$ denote real site $i$.
Using $\sum_{\mathbf{k}} \exp[i\mathbf{k}\cdot(\mathbf{R}_i-\mathbf{R}_j)] = N\delta_{ij}$, the inversion performed as
\begin{align*}
    \sum_{\mathbf{k}} \psi_{\mathbf{k}}(\mathbf{r}) e ^{-i \mathbf{k} \mathbf{R}_j} & = \frac{1}{\sqrt{N}}\sum_{\mathbf{k} i} w(\mathbf{r}-\mathbf{R}_i)e^{i\mathbf{k}\cdot(\mathbf{R}_i-\mathbf{R}_j)} \\
    \sum_{\mathbf{k}} \psi_{\mathbf{k}}(\mathbf{r}) e ^{-i \mathbf{k} \mathbf{R}_j} & = \sqrt{N} w(\mathbf{r}- \mathbf{R}_i)
\end{align*}
From this, we define
\begin{equation*}
    a_{i\sigma} = \frac{1}{\sqrt{N}}\sum_{\mathbf{k}} a_{\mathbf{k}\sigma} e^{i\mathbf{k}\cdot\mathbf{R}_i} \qquad
    a_{i\sigma}^{\dagger} = \frac{1}{\sqrt{N}}\sum_{\mathbf{k}} a_{\mathbf{k}\sigma}^{\dagger} e^{-i\mathbf{k}\cdot\mathbf{R}_i}
\end{equation*}
its inversion
\begin{equation*}
    a_{\mathbf{k}\sigma} = \frac{1}{\sqrt{N}}\sum_{i} a_{i \sigma}e^{-i\mathbf{k}\cdot\mathbf{R}_i} \\
    a_{\mathbf{k}\sigma}^{\dagger} = \frac{1}{\sqrt{N}}\sum_{i} a_{i \sigma}^{\dagger}e^{i\mathbf{k}\cdot\mathbf{R}_i}
\end{equation*}
and
\begin{equation*}
    \ket{\mathbf{k } } = \frac{1 }{\sqrt{N } } \sum_{i } \ket{w_i } e ^{i \mathbf{k }\mathbf{R_i}} \qquad
    \ket{w_i } = \frac{ 1 }{\sqrt{N }  } \sum_{\mathbf{k }} \ket{\mathbf{k} } e ^{-i \mathbf{k }\mathbf{R_i}}
\end{equation*}
so that
\begin{align*}
    \sum_{\mathbf{k} } \ket{\mathbf{k}} a_{\mathbf{k}\sigma}         & =  \frac{1}{\sqrt{N}}\sum_{\mathbf{k} i} \ket{\mathbf{k}}  a_{i \sigma}e^{-i\mathbf{k}\cdot\mathbf{R}_i} = \sum_{i}\ket{w_i} a_{i \sigma}                  \\
    \sum_{\mathbf{k} } \bra{\mathbf{k}} a_{\mathbf{k}\sigma}^\dagger & =  \frac{1}{\sqrt{N}}\sum_{\mathbf{k} i} \bra{\mathbf{k}}  a_{i \sigma}^\dagger e^{i\mathbf{k}\cdot\mathbf{R}_i} = \sum_{i}\bra{w_i} a_{i \sigma} ^\dagger
\end{align*}
The one particle operator transform as
\begin{equation*}
    H_1  =  \sum_{\mathbf{k} \sigma} \braket{ \mathbf{k }\sigma | h | \mathbf{k }\sigma} a^\dagger_{ \mathbf{k} \sigma} a_{ \mathbf{k} \sigma}
    =  \sum_{ij\sigma } \braket{w_i \sigma | h|w_j \sigma }a ^\dagger_{i \sigma} a_{j \sigma}
    =  \sum_{ij \sigma} \epsilon_{ij} a ^\dagger_{i \sigma} a_{j \sigma}
\end{equation*}
Whereas the two-particle operator
\begin{align*}
    H_2 & =  \frac{1}{2}\sum_{\substack{\mathbf{k}_1 \mathbf{k}_2 \sigma                                                        \\ \mathbf{k}_1' \mathbf{k}_2' \sigma'}}
    \braket{ \mathbf{k}_1   \mathbf{k}_2  | \frac{1}{r} |   \mathbf{k}_1' \mathbf{k}_2'   }
    a^\dagger_{ \mathbf{k}_1  \sigma} a^\dagger_{  \mathbf{k}_2  \sigma'} a_{\mathbf{k}_2'  \sigma'}a_{  \mathbf{k}_1'  \sigma} \\
    H_2 & =  \frac{1}{2}\sum_{\substack{i j \sigma                                                                              \\ k l \sigma'}}
    \braket{ i   j  | \frac{1}{r} |   k l   }
    a^\dagger_{ i  \sigma} a^\dagger_{  j  \sigma'} a_{l  \sigma'}a_{  k  \sigma}
\end{align*}

Thus, the Hamiltonian reads
\begin{equation*}
    H =  \sum_{ij \sigma} \epsilon_{ij} a ^\dagger_{i \sigma} a_{j \sigma}  +
    \frac{1}{2}\sum_{\substack{i j \sigma \\ k l \sigma'}}
    \braket{ i   j  | \frac{1}{r} |   k l   }
    a^\dagger_{ i  \sigma} a^\dagger_{  j  \sigma'} a_{l  \sigma'}a_{  k  \sigma}
\end{equation*}
with $\ket{w_i} = \ket{i} $ and the same goes for all indices.
The component can be split to shown how each contribute into the Hamiltonian
\begin{multline*}
    H = \sum_{ij\sigma} \epsilon_{ij}\,a_{i\sigma}^{\dagger}a_{j\sigma}
    + \frac{I}{2}\sum_{i\sigma}n_{i\sigma}n_{i-\sigma}
    + \frac{1}{2}\sum_{ij}^{ }{} '  \sum_{\sigma\sigma'}\langle ij|\frac{1}{r}|ij\rangle\,n_{i\sigma}n_{j\sigma'}\\
    -\frac{1}{2}\sum_{ij}{}' \sum_{\sigma\sigma'}J_{ij}a_{i\sigma}^{\dagger}a_{i\sigma'}a_{j\sigma'}^{\dagger}a_{j\sigma}
    + \sum_{ij} {} ' \sum_{\sigma\sigma'}\langle ii|\frac{1}{r}|jj\rangle\,a_{i\sigma}^{\dagger}a_{i\sigma'}^{\dagger}a_{j\sigma'}a_{j\sigma}
\end{multline*}

The second term occur at $i=j=k=l$.
The number particle operator appears after considering that the sum over spin having four terms, where each spin can take the value of up and down
\begin{equation*}
    \frac{1}{2}
    \sum_{i\sigma\sigma'} \braket{ii| \frac{1}{r} | ii } a^\dagger_{i\sigma} a^\dagger_{i\sigma'} a_{i\sigma'} a_{i\sigma} =
    \frac{I}{2}  \sum_{i,\sigma} a^\dagger_{i\sigma} a^\dagger_{i\sigma} a_{i\sigma} a_{i\sigma}
    +\frac{I}{2} \sum_{i,\sigma\neq\sigma'} a^\dagger_{i\sigma} a^\dagger_{i\sigma'} a_{i\sigma'} a_{i\sigma}
\end{equation*}
The matrix element reads
\begin{equation*}
    I =e^{2}\int d \mathbf{r}\;d \mathbf{r}'\;\frac{\big|w(\mathbf{r}-\mathbf{R}_{i})\big|^{2}\,\big|w(\mathbf{r}'-\mathbf{R}_{i})\big|^{2}}{|\mathbf{r}'-\mathbf{r}|}
    =e^{2} \int  d\mathbf r \;d\mathbf r'\; \frac{|w(\mathbf r)|^{2}\,|w(\mathbf r')|^{2}}{|\mathbf r'-\mathbf r|}
\end{equation*}
It is the energy cost for putting two electrons on the same site.
In the last term, we make the substitution $\mathbf{r} \rightarrow \mathbf{r} - \mathbf{R}_i $ so the integral become independent of $i$.
Only $\sigma \neq \sigma'$ survives, so
\begin{equation*}
    \frac{I}{2} \sum_{\sigma\neq\sigma'} a^\dagger_{i\sigma} a^\dagger_{i\sigma'} a_{i\sigma'} a_{i\sigma} = \frac{I}{2} \sum_{\sigma} a^\dagger_{i\sigma}a_{i\sigma} a^\dagger_{i-\sigma} a_{i-\sigma}
\end{equation*}

The third term occur at $i=k$ and $k=l$.
\begin{equation*}
    \frac{1}{2}\sum_{ij}{}' \sum_{\sigma \sigma'}
    \braket{ i   j  | \frac{1}{r} |   ij   }
    a^\dagger_{ i  \sigma} a^\dagger_{  j  \sigma'} a_{j  \sigma'}a_{i \sigma}
    = \frac{1}{2}\sum_{ij}^{ }{} '  \sum_{\sigma\sigma'}\langle ij|\frac{1}{r}|ij\rangle\,n_{i\sigma}n_{j\sigma'}
\end{equation*}

The fourth term occur at $i=l$ and $j=k$.
\begin{equation*}
    \frac{1}{2}\sum_{ij}{}' \sum_{\sigma \sigma'}
    \braket{ i   j  | \frac{1}{r} | ji   }
    a^\dagger_{ i  \sigma} a^\dagger_{  j  \sigma'} a_{i  \sigma'}a_{  j  \sigma}
    = -\frac{1}{2}\sum_{ij}{}' \sum_{\sigma\sigma'}J_{ij}a_{i\sigma}^{\dagger}a_{i\sigma'}a_{j\sigma'}^{\dagger}a_{j\sigma}
\end{equation*}
with the exchange integral of
\begin{equation*}
    J_{ij} = e^{2}\int d \mathbf{r}\;d \mathbf{r}' \frac{w ^*(\mathbf{r}-\mathbf{R}_{i})w(\mathbf{r}-\mathbf{R}_{j})w ^*(\mathbf{r}'-\mathbf{R}_{j})w(\mathbf{r}'-\mathbf{R}_{i})}{|\mathbf{r}'-\mathbf{r}|}
\end{equation*}

The fifth term occur at $i=j$ and $l=k$, with no further algebraic manipulation performed.

As the energy to put two electrons in single state approaching infinite $I \to \infty$, the state is single occupied.
This state is given by Block-Zener state
\begin{equation*}
    \lvert \sigma_1 \dots \sigma_N \rangle  a^\dagger_{1\sigma(1)} a^\dagger_{2\sigma(2)}\dots a^\dagger_{N\sigma(N)}\lvert 0\rangle
\end{equation*}
with $\sigma(i)$ specifies whether the electron at site $i$ is spin up or down.
The Hamiltonian
\begin{multline*}
    H = \sum_{ij\sigma} \epsilon_{ij}\,a_{i\sigma}^{\dagger}a_{j\sigma}
    + \frac{I}{2}\sum_{i\sigma}n_{i\sigma}n_{i-\sigma}
    + \frac{1}{2}\sum_{ij}^{ }{} '  \sum_{\sigma\sigma'}\langle ij|\frac{1}{r}|ij\rangle\,n_{i\sigma}n_{j\sigma'}\\
    -\frac{1}{2}\sum_{ij}{}' \sum_{\sigma\sigma'}J_{ij}a_{i\sigma}^{\dagger}a_{i\sigma'}a_{j\sigma'}^{\dagger}a_{j\sigma}
    + \sum_{ij} {} ' \sum_{\sigma\sigma'}\langle ii|\frac{1}{r}|jj\rangle\,a_{i\sigma}^{\dagger}a_{i\sigma'}^{\dagger}a_{j\sigma'}a_{j\sigma}
\end{multline*}
has the matrix element in this state
\begin{multline*}
    \braket{\left\{ \sigma_i(i) \right\}  |H |\left\{ \sigma_i(i)  \right\} } = \prod_{i=1}^N \delta_{\sigma'(i)\sigma(i)} \left[N\varepsilon_{ii}+\frac{1}{2}\sum_{ij} {}' \langle ij|\frac{1}{r}|ij\rangle\right] \\
    -\frac{1}{2}\sum_{ij} {}' \sum_{\sigma\sigma'}J_{ij}
    \braket{\left\{ \sigma_i(i) \right\}  |a^\dagger_{i\sigma}a_{i\sigma}a^\dagger_{j\sigma'}a_{j\sigma'}|\left\{ \sigma_i(i) \right\} }
\end{multline*}
The first three terms contribute diagonal component that is independent of spin.
As $I \rightarrow \infty$, $J_{ij}$ determine the state of the system.
The next step is then to determine the effective Hamiltonian in terms of this spin operator.

Based on their action on up and down  state, the product of annihilation and creation can be written as
\begin{equation*}
    \begin{matrix}
        a_{\uparrow}^{\dagger} a_{\uparrow} \equiv \begin{pmatrix}1 & 0 \\ 0 & 0\end{pmatrix}
         &
        a_{\uparrow}^{\dagger} a_{\downarrow} \equiv \begin{pmatrix}0 & 1 \\ 0 & 0\end{pmatrix}
        \\\\
        a_{\downarrow}^{\dagger} a_{\uparrow} \equiv \begin{pmatrix}0 & 0 \\ 1 & 0\end{pmatrix}
         &
        a_{\downarrow}^{\dagger} a_{\downarrow} \equiv \begin{pmatrix}0 & 0 \\ 0 & 1\end{pmatrix}
    \end{matrix}
\end{equation*}
Or in terms of Pauli matrices
\begin{align*}
    a_{\uparrow}^{\dagger} a_{\uparrow}   & = \frac{1 }{2 }\left[ \begin{pmatrix}    1&0\\1&-1   \end{pmatrix} +1\right] = \frac{1 }{2 }\left[ \sigma_3+1 \right]                                                     \\
    a_{\uparrow}^{\dagger} a_{\downarrow} & = \frac{1 }{2 }\left[ \begin{pmatrix}    0&1\\1&0   \end{pmatrix} +i \begin{pmatrix}    0&-i\\i&0   \end{pmatrix}\right] = \frac{1 }{2 }\left[ \sigma_1+i\sigma_y \right] \\
    a_{\downarrow}^{\dagger} a_{\uparrow} & = \frac{1 }{2 }\left[ \begin{pmatrix}    0&1\\1&0   \end{pmatrix} -i \begin{pmatrix}    0&-i\\i&0   \end{pmatrix}\right] = \frac{1 }{2 }\left[ \sigma_1-i\sigma_y \right] \\
    a_{\uparrow}^{\dagger} a_{\uparrow}   & = \frac{1 }{2 }\left[ 1 - \begin{pmatrix}    1&0\\1&-1   \end{pmatrix}\right] = \frac{1 }{2 }\left[ 1 -\sigma_3 \right]
\end{align*}
Next we written the operator in exchange term
\begin{align*}
    \sum_{\sigma\sigma'} a_{i\sigma}^\dagger a_{i\sigma'} a_{j\sigma'}^\dagger a_{j\sigma} 
    & =   a_{i\uparrow}^\dagger a_{i\uparrow} a_{j\uparrow}^\dagger a_{j\uparrow} + a_{i\uparrow}^\dagger a_{i\downarrow} a_{j\downarrow}^\dagger a_{j\uparrow}  \\
    &\qquad+ a_{i\downarrow}^\dagger a_{i\downarrow} a_{j\downarrow}^\dagger a_{j\downarrow}    + a_{i\downarrow}^\dagger a_{i\uparrow} a_{j\uparrow}^\dagger a_{j\downarrow}\\
    & = \frac{1}{4}\big[(\sigma_{3i}+1)(\sigma_{3j}+1)+(\sigma_{1i}+i\sigma_{2i})(\sigma_{1j}-i\sigma_{2j})\\
    &\qquad+(\sigma_{1i}-i\sigma_{2i})(\sigma_{1j}+i\sigma_{2j})+(1-\sigma_{3i})(1-\sigma_{3j})\big]\\
    & = \frac{1 }{2 }\left( \sigma_{1i  }\sigma_{1j} +\sigma_{2i  }\sigma_{2j}+\sigma_{3i  }\sigma_{3j}+1\right) = \frac{1 }{2 } (\boldsymbol{\sigma}_i \cdot \boldsymbol{\sigma}_j+1)
\end{align*}
Thus, neglecting the diagonal term, which are immaterial in determining the spin configuration, we have the Heisenberg Hamiltonian
\begin{equation*}
    H_H = -\frac{1 }{4 }\sum_{ij}{}' J_{ij} \boldsymbol{\sigma}_i \cdot \boldsymbol{\sigma}_j  
\end{equation*}
\end{document}